\documentclass{book}
\usepackage[utf8]{inputenc}
\usepackage[T1]{fontenc}
\usepackage{geometry}
\usepackage{graphicx}
\usepackage{float}
\usepackage{hyperref}
\usepackage[english,indonesian]{babel}

\title{Optimasi Pemrosesan Template LaTeX Menggunakan Pendekatan Berbasis AI}
\author{Academic Support Agent}
\date{\today}

\begin{document}
\maketitle

\begin{abstract}

\end{abstract}

\chapter{PENDAHULUAN}


\chapter{TINJAUAN PUSTAKA}


% Original content (commented out):
% \documentclass[a4paper,12pt,oneside]{book}

%Layout de página
\usepackage{geometry}\geometry{a4paper,total={170mm,257mm},left=20mm,top=20mm,}
\usepackage[font=it]{caption}
\usepackage{paralist}
\usepackage{multicol} %Para criar texto em colunas.

%Pacotes utilizados
\usepackage[utf8]{inputenc}
\usepackage[portuguese]{babel}
\usepackage[T1]{fontenc}
\usepackage{amsmath}
\usepackage{amsfonts}
\usepackage{amssymb}
\usepackage{graphics}
\usepackage{float}
\usepackage{subfigure}
\usepackage{xcolor}
\usepackage{array}
\usepackage{verbatim}
\usepackage{indentfirst}
\usepackage{tikz} %Para desenhar (ex. circuitos).
\usepackage{framed} % Moldura no texto
\usepackage{fancybox} % Para molduras adicionais nas caixas
\usepackage{textcomp}
\usepackage{circuitikz} %Desenhar circuitos elétricos.
\usepackage{wrapfig}
\usepackage{comment}
\usepackage{filecontents} % Clickable URLs (optional)
\usepackage[colorlinks=true,linktocpage,bookmarks=true,citecolor=black,linkcolor=black,urlcolor=black,hyperfootnotes=true,pageanchor=true]{hyperref} %Clicáveis em cor preta

%Configurações de fontes e tipografia
\linespread{1.5} %Espaçamento entre as linhas do texto
\usepackage{helvet} %Fonte do texto (Arial)
\renewcommand{\familydefault}{\sfdefault}
\usepackage{graphicx}
\usepackage{ragged2e} %Texto Justificado

%Pacotes para tabelas.
\usepackage{multirow} 
\usepackage{tabularx}
\usepackage{scalefnt} % Redimensioar tabelas
\usepackage{chngcntr} %Enumerar as tabelas e figs de acordo com o capítulo.
\counterwithin{figure}{chapter}
\counterwithin{table}{chapter}

%Bibliografia
\usepackage{lipsum}
\usepackage{natbib}
\usepackage{cite}


%================================================================================================================
%==========================================INÍCIO DO DOCUMENTO===================================================
%================================================================================================================

\begin{document}

%----------------------------------------------------------------------------------------------------------------
%Logo da Instituição - MNPEF, UnB e SBF
\begin{figure}[H]
	\subfigure{\includegraphics[width=3.4cm]{Imagens/Logo_MNPEF.png}\hspace*{2.5cm} \includegraphics[width=3.6cm]{Imagens/Logo_UnB.png}\hspace*{2.5cm} \includegraphics[width=4.5cm]{Imagens/Logo_SBF.jpg}}
\end{figure}

%----------------------------------------------------------------------------------------------------------------
%------------------------------------------------CAPA 1 - Título e Autor-----------------------------------------
%----------------------------------------------------------------------------------------------------------------
\vspace{3cm}	

\begin{center}	
TÍTULO DA DISSERTAÇÃO
\end{center}

\vspace{3cm}	

\begin{center}
Nome do Mestrando
\end{center}

\vspace{3cm}

\hfill \parbox{8.5cm} {\noindent Dissertação de Mestrado apresentada ao Programa de Pós-Graduação (nome dado na instituição) no Curso de Mestrado Profissional de Ensino de Física (MNPEF), como parte dos requisitos necessários à obtenção do título de Mestre em Ensino de Física.}
\vfill

\vspace{1.0cm}

\hfill \parbox{8.5cm} {\noindent Orientador(es):
	
	Nome do Orientador 1
	
	Nome do Orientador 2}
\vfill

\vspace{2.0cm}

\begin{center}
Local

Mês e Ano da Defesa
\end{center}

\thispagestyle{empty} %Remove a numeração desta página
\newpage
%----------------------------------------------------------------------------------------------------------------
%-------------------------------------------------%CAPA 2 DA DISSERTAÇÃO-----------------------------------------
%----------------------------------------------------------------------------------------------------------------
\begin{center}	
	TÍTULO DA DISSERTAÇÃO
\end{center}

\vspace{0.5cm}	

\begin{center}
	Nome do Mestrando
\end{center}

\vspace{0.5cm}

\begin{center}
Orientador(es):

Nome do Orientador 1

Nome do Orientador 2  
\end{center} 

\vspace{1.0cm}

\begin{center} \justifying
\noindent Dissertação de Mestrado submetida ao Programa de Pós-Graduação (nome dado na instituição) no Curso de Mestrado Profissional de Ensino de Física (MNPEF), como parte dos requisitos necessários à obtenção do título de Mestre em Ensino de Física.
\end{center}

\vspace{0.5cm}

\noindent Aprovada por

\vspace{2cm}

%Banca Examinadora
\hfill \parbox{8.5cm} {\rule{9cm}{1pt} %Linha com dimensões
	
Dr. Nome do Membro da Banca

\vspace{1.0cm}

\rule{9cm}{1pt} %Linha com dimensões

Dr. Nome do Membro da Banca

\vspace{1.0cm}

\rule{9cm}{1pt} %Linha com dimensões

Dr. Nome do Membro da Banca}
\vfill

\vspace{2.0cm}

\begin{center}
	Local
	
	Mês e Ano da Defesa
\end{center}

\pagestyle{plain}
\pagenumbering{roman}
\setcounter{page}{2} %Escolhe-se o estilo de numeração da página.
\newpage
%----------------------------------------------------------------------------------------------------------------
%-----------------------------------------PÁGINA DA FICHA CATALOGRÁFICA------------------------------------------
%----------------------------------------------------------------------------------------------------------------
\begin{center}
\textbf{FICHA CATALOGRÁFICA (MODELO)}
\end{center}

\vspace{1.0cm}

\begin{center}
\fbox{
	\begin{minipage}{0.9\textwidth} %largura de texto
		S586p
		
		Silva, Sergio Tobias da
		
		Propagação do Som: Conceitos e Experimentos / Sergio
		Tobias da Silva - Rio de Janeiro: UFRJ / IF, 2011.
		
		viii, 77 f.: il.;30cm.
		
		Orientador: Carlos Eduardo Aguiar
		
		Dissertação (mestrado) – UFRJ / Instituto de Física /		
				Programa de Pós-Graduação em Ensino de Física, 2011.
				
		Referências Bibliográficas: f. 74-77.
		
		1. Ensino de Física. 2. Som. 3. Velocidade do som. I. Aguiar,
		Carlos Eduardo. II. Universidade Federal do Rio de Janeiro, Instituto
		de Física, Programa de Pós-Graduação em Ensino de Física. III.		
		Propagação do Som: Conceitos e Experimentos.
	\end{minipage}}
\end{center}
\newpage
%----------------------------------------------------------------------------------------------------------------
%------------------------------------------PÁGINA DA CITAÇÃO OU FRASE--------------------------------------------
%----------------------------------------------------------------------------------------------------------------

\null
\vfill

\hfill \parbox{9.0cm} {\noindent Ora, a realidade é constituída por essências e existências particulares e, portanto, o conhecimento verdadeiro tem que ser um conhecimento que preserve o particular sem destruí-lo numa nomenclatura abstrata. 

\vspace{0.5cm}

\noindent Baruch de Espinosa}

\newpage
%----------------------------------------------------------------------------------------------------------------
%----------------------------------------PÁGINA DE DEDICATÓRIA---------------------------------------------------
%----------------------------------------------------------------------------------------------------------------

\null
\vfill

\hfill \parbox{10.0cm} {\noindent Dedicado a minha família... etc, etc e etc.}

\newpage
%----------------------------------------------------------------------------------------------------------------
%----------------------------------------PÁGINA DE AGRADECIMENTOS------------------------------------------------
%----------------------------------------------------------------------------------------------------------------
\begin{center}
\textbf{Agradecimentos}
\end{center}
	
Se for o caso agradecer à CAPES pelo apoio financeiro por meio da bolsa
concedida.

\newpage
%----------------------------------------------------------------------------------------------------------------
%----------------------------------------------------RESUMO------------------------------------------------------
%----------------------------------------------------------------------------------------------------------------

\begin{center}
	\textbf{RESUMO}
\end{center}

\begin{center}	
	TÍTULO DA DISSERTAÇÃO
\end{center}

\begin{center}
	Nome do Mestrando
\end{center}

\begin{center}
	Orientador(es):
	
Nome do Orientador 1

Nome do Orientador 2  
\end{center} 

\begin{center} \justifying
	\noindent Dissertação de Mestrado submetida ao Programa de Pós-Graduação (nome dado na instituição) no Curso de Mestrado Profissional de Ensino de Física (MNPEF), como parte dos requisitos necessários à obtenção do título de Mestre em Ensino de Física.
\end{center}


\noindent Texto do resumo...


Palavras-chave: Ensino de Física, palavra chave 2, palavra chave 3.

\vspace{14cm} %Alterar a distancia quando for digitar o resumo.

\begin{center}
Local

Data
\end{center}

\newpage
%----------------------------------------------------------------------------------------------------------------
%--------------------------------------------------ABSTRACT------------------------------------------------------
%----------------------------------------------------------------------------------------------------------------

\begin{center}
	\textbf{ABSTRACT}
\end{center}

\begin{center}	
	TITLE OF DISSERTATION
\end{center}

\begin{center}
	Nome do Mestrando
\end{center}

\begin{center}
	Supervisor(s):
	
Nome do Orientador 1

Nome do Orientador 2 
\end{center} 

\begin{center} \justifying
	\noindent Abstract of master’s thesis submitted to Programa de Pós-Graduação (nome dado na instituição) no Curso de Mestrado Profissional de Ensino de Física (MNPEF), in partial fulfillment of the requirements for the degree Mestre em Ensino de Física.
\end{center}


\noindent Abstract text...


Keywords: Physics education, keyword 2, keyword 3

\vspace{14cm} %Alterar a distancia quando for digitar o resumo.

\begin{center}
Local

Month and Year
\end{center}

\newpage
%----------------------------------------------------------------------------------------------------------------
%--------------------------------------------------LISTA DE FIGURAS----------------------------------------------
%----------------------------------------------------------------------------------------------------------------

\listoffigures

\newpage
%----------------------------------------------------------------------------------------------------------------
%-----------------------------------------------LISTA DE TABELAS-------------------------------------------------
%----------------------------------------------------------------------------------------------------------------

\listoftables

\newpage
%================================================================================================================
%===================================================SUMÁRIO======================================================
%================================================================================================================

\renewcommand{\contentsname}{Sumário}

\tableofcontents
%\addcontentsline{entrada}

\newpage
%================================================================================================================
%=============================================INÍCIO DO TRABALHO=================================================
%================================================================================================================

\chapter{Título do Capítulo 1}

\section{Título da seção}

A dissertação deverá ser apresentada em papel branco, tamanho A4, com
margens de aproximadamente 3 cm. A fonte utilizada no corpo principal do texto não
deve ter tamanho inferior a 10 pt ou superior a 12 pt, salvo em partes especiais como
títulos, legendas, fórmulas, notas de rodapé, sobrescritos, etc. O tamanho recomendado
é 12 pt. Essa é a fonte utilizada no presente texto.

O espaço entre as linhas do corpo principal do texto não deve ser inferior a 1 ou
superior a 1,5, salvo em partes especiais como títulos, legendas, fórmulas, notas de
rodapé, etc. O espaçamento recomendado é 1,5, semelhante ao utilizado neste texto. Os
parágrafos do corpo principal do texto devem ter alinhamento justificado.
Todas as páginas da dissertação, com exceção da folha de rosto e da que começa
a parte textual, deverão ser numeradas. São usados algarismos romanos (minúsculos)
para as páginas pré-textuais e algarismos arábicos nas páginas textuais e pós-textuais.


\pagestyle{plain}
\pagenumbering{arabic}
\setcounter{page}{1} %Escolhe-se o estilo de numeração da página.
\newpage

%Capítulo 2
%----------------------------------------------------------------------------------------------------------------
\chapter{Título do capítulo 2}

\section{Título da seção}

Os capítulos da dissertação devem ter título e ser numerados sequencialmente,
em algarismos arábicos ou romanos. Capítulos podem ser divididos em seções, que
devem ter título e ser numeradas sequencialmente em algarismos arábicos, na forma
“capítulo. seção” (por exemplo, 2.3 ou II.3).
Seções podem ser divididas em subseções, que devem ter título e ser numeradas
sequencialmente em algarismos arábicos, na forma “capítulo.seção.subseção” (por
exemplo, 2.3.1 ou II.3.1).
Notas de rodapé são identificadas por numeração sequencial, em algarismos
arábicos, como se vê aqui\footnote{Exemplo de nota de rodapé.}. Elas devem ser referenciadas com o número sobrescrito e
são impressas ao pé da página, separadas do texto normal por uma barra horizontal.
Recomenda-se adotar fonte menor que a usada no corpo do texto. As notas devem ser
colocadas na página em que aparecem as chamadas numéricas, evitando a continuação
nas páginas seguintes.

\subsubsection{Título da subseção}

Citações diretas, consistindo na transcrição literal de palavras ou trechos de
outro autor, podem ser apresentadas de duas formas:

\begin{itemize}

\item[i.] Inseridas entre aspas, no meio do texto normal, se ocuparem cerca de três linhas;
 
\item[ii.] Impressas em destaque na folha (com recuo na margem esquerda, ou em espaço
menor), no caso de citações mais longas, também entre aspas.

\end{itemize}

Citações diretas devem vir acompanhadas da identificação do autor citado ou de
menção à referência bibliográfica correspondente. A reprodução de figura ou elemento
gráfico de outra obra deve ser autorizada por escrito pelo detentor dos direitos autorais.
O texto da dissertação deve fazer menção à autorização e citar a obra original. Mesmo
no caso de material de domínio público, a citação da obra original é obrigatória.

\subsubsection{Título da subseção}

As referências bibliográficas devem conter as seguintes informações:

\begin{itemize}

\item[i.] \textit{Artigo em periódico}: autor, título do artigo, nome da revista, número do volume,
número do fascículo, páginas ou identificação do artigo, ano de edição.

\item[ii.] \textit{Artigo em anais de conferência}: autor, título do artigo, título dos anais, nome,
local e ano da conferência, local e nome da editora, ano de edição, páginas ou
identificação do artigo.

\item[iii.] \textit{Livro}: autores, título do livro, número da edição, local e nome da editora, data da
edição, páginas selecionadas.

\item[iv.] \textit{Capítulo de livro}: autores do capítulo, título do capítulo, título do livro,
organizadores do livro, local e nome da editora, data da edição, páginas
selecionadas.

\item[v.] \textit{Tese ou dissertação}: autor, título da obra, nível (mestrado ou doutorado),
instituição onde foi apresentada, ano da defesa, páginas selecionadas.

\item[vi.] \textit{Página da internet}: autor ou instituição ou companhia responsável pela página,
informações complementares (se houver), endereço eletrônico (url), data de
acesso.

\end{itemize}

Citações a referências bibliográficas podem ser feitas numericamente, em ordem de apresentação ([1], [2], etc.), ou com o sobrenome do primeiro autor seguido do ano de publicação (por exemplo, [Silva 2003]), ou ainda abreviando a última forma (por exemplo, [Sil03]). No caso de páginas da internet, podem ser usados uma identificação da página e o ano de acesso (por exemplo, [Wiki-Termo 2010] para referências ao conteúdo de \url{http://pt.wikipedia.org/wiki/Termodinâmica} em 2010). As entradas das referências bibliográficas devem estar padronizadas com o formato de citação no corpo do trabalho. A lista de referências bibliográficas deve ser ordenada de acordo com suas entradas, em sequência numérica ou alfabética (dependendo do formato de citação escolhido).

\subsection{Título da seção}

Figuras e tabelas devem ter numeração sequencial, em algarismos arábicos. A seqüência pode ser feita, opcionalmente, dentro de cada capítulo. Nesse caso a numeração terá a forma “capítulo.número” (por exemplo, 4.3 ou IV.3). A numeração é parte da legenda, que deve vir abaixo da figura ou tabela, mas sempre na mesma página. Exemplos de formatação de legendas estão mostrados a seguir.

%Tabela 2.1
\begin{table}[H]
 \begin{center}
 \begin{tabular}{ |p{4cm}|p{10cm}| }
\hline
\begin{center} \textbf{Aspecto conceitual} \end{center}   & \begin{center}  \textbf{Concepções intuitivas} \end{center}   \\ \hline
    A transmissão do som      &  O som é uma propriedade da fonte de som, algo que não é transmitido.
                                 
                                 O som é algo material, um objeto que se move de um lugar para outro.
                                 
                                 O som é uma substância discreta, “uma coleção única de partículas ou moléculas”     que é transportada de um local a outro (por ex. o ar e o vento).
                                 
                                 O som é transmitido na água porque há oxigênio ou ar na água.
                                 
                                 O som é transmitido nos sólidos através de pequenos buracos ou aberturas no interior do material. \\ \hline
    
 A velocidade do som           &  Som de baixa frequência (som grave, abafado) move-se a uma velocidade diferente do som de alta freqüência (estridente).
    
    Sons de alto volume (sons fortes) são transmitidos mais rapidamente que os sons de baixo volume (sons fracos). \\ \hline
  \end{tabular}
  \caption{Exemplo de tabela com legenda.\label{Tabela_2.1}}
\end{center}
\end{table}

%Inserção de Imagem - 1
\begin{figure}[H]
\centering
\includegraphics[scale=0.27]{Imagens/Referenciais_Galileu.png}
\caption{Exemplo de figura com legenda. \label{Fig_2.1}}
\end{figure}

\newpage
%================================================================================================================
%=============================================APÊNDICES==========================================================
%================================================================================================================

\appendix

\chapter{Título do apêndice}

Apêndices ou anexos devem ter título e ser identificados sequencialmente por
letras (A, B, etc.) ou por números em algarismos arábicos ou romanos. O produto
educacional que acompanha a dissertação de mestrado deve, preferencialmente, ser
apresentado em destaque em um ou mais apêndices. Anexos diferenciam-se de
apêndices por serem de autoria de terceiros.


\chapter{Título do apêndice}

Este será o texto do APÊNDICE B.


\newpage
%================================================================================================================
%=====================================REFERÊNCIAS BIBLIOGRÁFICAS=================================================
%================================================================================================================

%\renewcommand{\bibname}{Referências Bibliográficas}

\phantomsection
\addcontentsline{toc}{chapter}{Bibliografia}

\nocite{caruso2006fisica} %Livro
\nocite{secretaria1999} %Cap. de livro.
\nocite{costa1998} %Artigo.
\nocite{birds1998} %Site
\nocite{gomes1998} %Livro

\bibliographystyle{plain} % Define o estilo das referências
\bibliography{Modelo_de_referencias}  %Nome(s) do(s) arquivo(s) que contém as referências arquivo na pasta do .tex:(Modelo_de_referencias).


%====================================================FIM=========================================================
\end{document}

%Daniel R. Ferreira (MNPEF-UnB) - daniel.rf_ufg@hotmail.com - Modelo de dissertação em LateX, cortesia para o MNPEF.




\end{document}